\subsection{约化范数与行列式}

在本小节中，D 是 K 上的一个有限次数的体，其中心为 K。我们用 \(D_{ab}^{*}\) 表示乘法群 \(D^{*}\) 对其导群（换位子群）的商群，用 \(\pi\) 表示从 \(D^{*}\) 到 \(D_{ab}^{*}\) 的典范同态。映射 \(Nrd_{D/K}\) 诱导一个从 \(D^{*}\) 到 \(K^{*}\) 的群同态；由于 K 是交换的，这个同态的核包含 \(D^{*}\) 的导群。因此，存在唯一的同态 \(\overline{Nrd}\)（从 \(D_{ab}^{*}\) 到 \(K^{*}\) 的），使得 \(\mathrm{Nrd}_{\mathrm{D/K}}(x)=\overline{\mathrm{Nrd}}\circ\pi(x)\) 对每个 \(x\in D^{*}\) 成立。

命题 9. —— 设 V 是体 D 上的一个有限维右向量空间。设 E 是体 K 上的代数 \(\operatorname{End}_{\mathrm{D}}(\mathrm{V})\)；它是中心单的且为有限次数。对 E 的每个可逆元 u，我们有

\[
\mathrm{Nrd} _ {\mathrm{E/K}} (u) = \overline {{\mathrm{Nrd}}} (\det u)\tag{64}
\]

（\(\det u\) 表示 \(u\) 的行列式；其定义参看 VIII, 第 452 页，命题 2）。

我们用 n 表示 V 在 D 上的维数，并利用 V 在 D 上的一组基把 E 与矩阵代数 \(\mathbf{M}_{n}(\mathrm{D})\) 等同起来。代数 E 的乘法群 \(\mathrm{GL}_{n}(\mathrm{D})\) 由对角矩阵与矩阵 \(\mathrm{B}_{ij}(\lambda)\) 生成（II, §10, No.~13, 第 362 页，推论 1）。因此，命题 9 由下面两个特殊情形得出。

\begin{enumerate}
\def\labelenumi{\Alph{enumi})}
\tightlist
\item
  假设 u 是对角矩阵 \(\mathrm{diag}(a_{1},\ldots,a_{n})\)。对每个 \(1\leqslant i\leqslant n\)，设 \(L_{i}\) 是 D 的一个包含 \(a_{i}\) 的极大交换子域；设 L 是 E 的由形如 \(\mathrm{diag}(t_{1},\ldots,t_{n})\)（其中 \(t_{i}\in L_{i}\)，\(1\leqslant i\leqslant n\)）的对角矩阵所组成的子代数。设 d 是 D 在 K 上的约化次数。我们有 \([L_{i}:K]=d\)（对 \(1\leqslant i\leqslant n\)；VIII, 第 265 页，推论 2）。K-代数 L 同构于 \(L_{1}\times\cdots\times L_{n}\)，因而是次数为 nd 的半单代数；而我们又有 \([E:K]=n^{2}[D:K]=n^{2}d^{2}=[L:K]^{2}\)。由此推出 L 是 E 的一个极大交换半单子代数（VIII, 第 262 页，命题 3）。由 VIII, 第 346 页，命题 7，我们有 \(\mathrm{Nrd}_{\mathrm{E}/\mathrm{K}}(u)=\mathrm{N}_{\mathrm{L}/\mathrm{K}}(u)\)。于是，由 III, §9, No.~3, 第 544 页的公式 (18)，我们有
\end{enumerate}

\[
\begin{array}{c} \mathrm{Nrd} _ {\mathrm{E/K}} (u) = \prod_ {i = 1} ^ {n} \mathrm{N} _ {\mathrm{L} _ {i} / \mathrm{K}} (a _ {i}) = \prod_ {i = 1} ^ {n} \mathrm{Nrd} _ {\mathrm{D/K}} (a _ {i}) \\ = \mathrm{Nrd} _ {\mathrm{D/K}} (a _ {1} \dots a _ {n}) = \overline {{\mathrm{Nrd}}} (\pi (a _ {1} \dots a _ {n}))  . \end{array}
\]

此外，由 VIII, 第 453 页，命题 3，我们有 \(\det u = \pi(a_{1} \cdots a_{n})\)，这在此情形给出了公式 (64)。

\begin{enumerate}
\def\labelenumi{\Alph{enumi})}
\setcounter{enumi}{1}
\tightlist
\item
  假设 u 等于 \(\mathrm{B}_{ij}(\lambda)\)，其中 \(\lambda\) 是 D 的一个元素，i, j 是区间 \([1,n]\) 中两个不同的整数。用 d 表示 D 在 K 上的约化次数，用 M 表示由 V 经标量从 D 到 K 的限制而得到的 K 上的向量空间。那么 M 是一个单 E-模，且我们有
\end{enumerate}

\[
\mathrm{Pc} _ {\mathrm{M/K}} (u; \mathrm{X}) = \mathrm{Pcrd} _ {\mathrm{E/K}} (u; \mathrm{X}) ^ {d}\tag{65}
\]

这由 VIII, 第 341 页的公式 (26) 得出。此外，M 是 K 上维数为 \(nd^2\) 的向量空间，且 \(u - 1_{\mathrm{M}}\) 是 M 的一个幂零自同态；因此我们有

\[
\mathrm{Pc} _ {\mathrm{M/K}} (u; \mathrm{X}) = (\mathrm{X} - 1) ^ {n d ^ {2}}.\tag{66}
\]

比较公式 (65) 与 (66)，我们得到

\[
\mathrm{Pcrd} _ {\mathrm{E/K}} (u; \mathrm{X}) = (\mathrm{X} - 1) ^ {n d}\tag{67}
\]

特别地，\(\mathrm{Nrd}_{\mathrm{E}/\mathrm{K}}(u)=1\)。我们又由 VIII, 第 453 页，命题 3 有 det u=1；因此公式 (64) 在此情形成立。

注. —— 若 E 的元素 u 不可逆，则我们有 \(\mathrm{Nrd}_{\mathrm{E}/\mathrm{K}}(u)=0\)（VIII, 第 340 页，命题 3）。

